\input{\string~/Documents/School/"UC Irvine"/ta_template2.0.tex}
\rh{2B}{11.8}
\chead{\today}
\graphicspath{ {\string~/Desktop} }
\usetikzlibrary{arrows}
%============ANSWER KEY TOGGLE===========
\newcommand{\ans}[1]{ \noindent {\color{red} \textbf{Answer:} #1}}
\renewcommand{\ans}[1]{\vspace{3in}}
%==========================================
\begin{document}
\begin{center}
\textbf{Power Series}
\end{center}

\begin{defi}[Power Series]
	A power series is a series of the form:
	\[
		\sum_{n =0}^\infty c_n x^n = c_0 + c_1 x + c_2 x^2 + \cdots 
	\]
	such that, for each value of $x$, we can test for convergence or divergence. More generally, the series:
	\[
		\sum_{n =0}^\infty c_n (x-a)^n = c_0 + c_1 (x-a) + c_2 (x-a)^2 + \cdots 
	\]
	is a power series centered at $x = a$. 
\end{defi}

\begin{thm}
	For a given power series $\sum_{n=0}^{\infty} c_n(x-a)^n$, there are only three possibilities:
	\begin{enumerate}[label=$(\roman*)$]
		\item The series converges only when $x=a$.
		\item The series converges for all $x$.
		\item There is a positive number $R$ such that the series converges if $|x-a|<R$ and diverges if $|x-a|>R$.
	\end{enumerate}
	In the instance of $(iii)$, we call $R$ the \emph{radius of convergence}.
\end{thm}

\begin{exx}
	Determine the radius of convergence of the following power series centered at $a = 0$:
	\[
		\sum_{n = 0}^\infty n! x^n
	\]
\end{exx}
\ans{Whenever we see a factorial, our heads should immediately think: \emph{ratio test}:
\[
\lim_{n \to \infty} \lrabs{\frac{a_{n+1}}{a_n}} = \lim_{n \to \infty} \lrabs{\frac{(n+1)!x^n}{n! x^n}} = \lim_{n \to \infty} \lrabs{(n + 1)x} < \infty
\]
where the last inequality holds if and only if $x = 0$. 
}

\begin{exx}
	Determine the radius of convergence and interval of convergence of the following power series. 
	\[
	\sum_{n = 0}^\infty \frac{n(x+2)^n}{3^{n+1}}
	\]
\end{exx}

\ans{Note first that this series is centered at $a = -2$, so our interval of convergence will require an additional step. We'll use the Ratio test:
\[
\lim_{n \to \infty} \lrabs{\frac{a_{n +1}}{a_n}} = \lim_{n \to \infty} \lrabs{\frac{n+1 (x + 2)^{n+1}}{3^{n+2}} \cdot \frac{3^{n+1}}{n(x+2)^n}} = \lim_{n \to \infty} \lrabs{\frac{n+1}{3n}(x+2)} = \frac{\lrabs{x+2}}{3}
\]
Using the conclusion of the Ratio test, we see that:
\[\sum_{n = 0}^\infty \frac{n(x+2)^n}{3^{n+1}}
\begin{cases}
	\text{converges}, &\text{ if } \frac{\lrabs{x+2}}{3} < 1\\
	\text{diverges}, &\text{ if } \frac{\lrabs{x+2}}{3} > 1\\
\end{cases}
\]
So we conclude that our $R = 3$ and the interval of convergence is $(-5, 1)$. As an exercise, test the points $-5, 1$ to see that the series doesn't converge at either, so we need not add either point to our interval of convergence. 
}


\problembox{
	Determine the radius of convergence and interval of convergence of the following power series. 
	\[
	\sum_{n=0}^{\infty} \frac{n+1}{(2 n+1) !}(x-2)^n
	\]
}

\ans{\href{https://tutorial.math.lamar.edu/Solutions/CalcII/PowerSeries/Prob3.aspx}{link}}

\problembox{
	Determine the radius of convergence and interval of convergence of the following power series. 
	\[
		\sum_{n=0}^{\infty} \frac{4^{1+2 n}}{5^{n+1}}(x+3)^n
	\]
}

\ans{\href{https://tutorial.math.lamar.edu/Solutions/CalcII/PowerSeries/Prob4.aspx}{Link}}
\end{document}